\section{Introduction}
\label{sec:intro}

The symmetries that organize low-energy physics---the Lorentz invariance of spacetime
and the $U(1)\times SU(2)\times SU(3)$ internal gauge group of the Standard Model---are
normally taken as primitive postulates. An alternative, long-running research program
asks whether some of them could instead be \emph{emergent}: low-energy, approximate
symmetries of a simpler underlying medium that carries no such symmetry fundamentally.
This viewpoint has a substantial pedigree. Emergent gauge fields and even emergent
relativistic fermions arise in superfluid ${}^3$He \citep{Volovik2003} and in
lattice bosonic models through string-net condensation \citep{Wen2004,LevinWen2005};
emergent Lorentzian ``acoustic'' metrics arise for phonons in moving media
\citep{Unruh1981,BarceloLiberatiVisser2005}; and gauge invariance itself can appear as
a low-energy attractor of a generic non-gauge-invariant lattice theory
\citep{ForsterNielsenNinomiya1980}. The mathematical engine common to many of these
constructions is the geometric phase: the abelian Berry phase
\citep{Berry1984,Simon1983} and its non-abelian Wilczek--Zee generalization
\citep{WilczekZee1984,ShapereWilczek1989}, whose curvature over a parameter manifold
behaves as a gauge field strength \citep{XiaoChangNiu2010}.

This paper studies whether a \emph{single, deterministic, purely elastic} substrate can
carry both an emergent Lorentzian light cone and an emergent $U(1)\times SU(3)$ gauge
structure through this geometric-phase mechanism. The substrate is minimal: a static
four-dimensional cubic lattice embedded with codimension zero in a 4D Euclidean ambient,
governed by a central-force spring action whose \emph{only} nonlinearity is the
Pythagorean (Euclidean) link length. No gauge field, no metric, and no internal symmetry
group is postulated; each is sought as a derived, effective object. The four ambient
directions are geometrically equivalent, and the distinction of one as ``time'' is a
property of the action's sign structure, not of the ambient.

\paragraph{What this paper establishes.} We report a set of explicit, reproducible
results, each emitted by a named numerical script (Appendix~\ref{app:repro}):
\begin{enumerate}
  \item \textbf{Lorentzian signature from prestress.} Linearizing the action yields an
  effective metric whose signature $(3,1)$ is exactly the sign pattern
  $\eta_\mu=(-1,-1,-1,+1)$ of the directional prestress; flipping the temporal sign
  gives a Euclidean, non-propagating theory. The branch speeds are given in closed form
  (Section~\ref{sec:metric}).
  \item \textbf{Full $\su{3}$ from a rank-3 carrier.} About a finite-amplitude periodic
  vacuum carrier, the Bloch fluctuation operator has an isolated rank-3 complex
  eigenbundle whose gauge-invariant Wilczek--Zee curvature generates \emph{all eight}
  generators of $\su{3}$---not merely the $\mathfrak{so}(3)$ of a real frame. The result
  is generic across backgrounds, leading order in the carrier amplitude, and satisfies
  the non-abelian Bianchi identity (Section~\ref{sec:carrier}).
  \item \textbf{Maxwell and Yang--Mills dynamics.} Emergent local gauge invariance,
  locality, and power counting force the leading effective action of the trace $U(1)$ and
  traceless $SU(3)$ sectors to be $-\tfrac14 f^2$ and $-\tfrac14 G^2$, with the mass term
  forbidden (massless photon and gluon) and substrate-induced, positive, finite couplings
  (Section~\ref{sec:gauge}). We further demonstrate that the same curvature appears over
  \emph{spacetime}, not only over the Brillouin zone, closing the interpretive gap
  between a $k$-space bundle and a physical gauge field.
  \item \textbf{Quantitative sector split.} The rank selection $U(3)$-not-$U(4)$ and the
  $U(1)/SU(3)$ split are controlled by the prestress parameters
  $(\alpha_s,\alpha_t,\gamma_t)$ (Section~\ref{sec:split}).
\end{enumerate}

\paragraph{What this paper does not claim.} We are deliberate about scope, and state the
principal limitations up front (they are quantified in Section~\ref{sec:limits}). First,
the specific color-carrying background is a stable but \emph{higher-energy texture}, not
the elastic ground state; so this is a demonstration that an $\su{3}$ sector can live on
a substrate texture, not that the ground-state vacuum is a Yang--Mills field. Second, the
emergent Lorentz invariance is only \emph{approximate}: the lattice is genuinely
anisotropic, and---as we show---this residual anisotropy is not a defect to be hidden but
the very property that lets the emergent massless gauge bosons evade the Weinberg--Witten
theorem \citep{WeinbergWitten1980}. Third, the matter sector (self-confined solitons,
mass, color sources, confinement) is deliberately deferred to future work; consequently
we make no contact with measured couplings or mass ratios. This is a constructive proof
of concept about emergent \emph{symmetry structure}, presented with its boundaries
explicit.

\paragraph{Novelty.} That holonomy can mimic gauge structure is well established; what is
new is that an abelian (Maxwell) and a non-abelian ($SU(3)$) sector, together with a
Lorentzian cone, descend from the \emph{same} carrier geometry of one deterministic
prestressed lattice, with the $\su{3}$ (versus $\mathfrak{so}(3)$) content exhibited
explicitly and every claim reproducible (Appendix~\ref{app:repro}).

\paragraph{Determinism and quantum foundations.} The substrate is fully deterministic
and ontic; probabilistic quantum behavior is not part of the present work. We note only,
and defer to future work, that a deterministic locally-mediated substrate is compatible
with Bell's theorem \citep{Bell1964} provided measurement independence is relaxed, which
for a time-symmetric variational (all-at-once) formulation is the ordinary consequence of
boundary data rather than a conspiracy---placing the model in the locally-mediated
time-symmetric class of Wharton and Argaman \citep{WhartonArgaman2020}. The Born rule is
not derived here and is treated as an open problem.
